This thesis proposes an end-to-end pipeline for generating 2.5D animated videos from a single
Traditional Chinese Painting (TCP) image.
Existing diffusion-based monocular depth estimation models tend to produce erroneous depth maps
for Chinese paintings due to domain gaps such as white-space semantics and abstract brushstrokes.
To address this, we incorporate Marigold-DC, enabling users to guide depth estimation with
sparse depth control points, and design \textit{ChinesePaintingLoss}, a composite loss function
tailored to Chinese painting characteristics, comprising brushstroke-aware smoothness loss,
Laplacian loss, whitespace regularization, and a ranking loss.
The refined depth map is then unprojected into a colored 3D point cloud using metric depth
estimated by Depth Anything V3, rasterized along a predefined camera trajectory to produce a
parallax-rich point cloud video, and finally fed into Uni3C's PCDController, where a video
diffusion model fills occlusion holes while preserving the painterly style, yielding a smooth,
complete 2.5D animated video.

\begin{flushleft}
\mbox{{\bf Keywords}: Chinese Painting, 2.5D Video, Depth Estimation, Diffusion Model, Point Cloud Rendering}
\end{flushleft}